\section{Discussion}
\label{sec:discussion}

\subsection{What a random encoder tells us}
The central practical proposal of this paper is methodological: fNIRS foundation-model papers
should report a randomly initialised encoder of the same architecture as a baseline. The control
is cheap---ten seeds cost minutes on a laptop CPU---and it changes the interpretation of a
result entirely. In our data the absolute accuracy of MtReNIRS ($0.469$) looks like a
respectable result for three-class fNIRS decoding on a sparse montage. Against a random encoder
($0.448 \pm 0.010$) it is a $2$-pp difference that 17 folds cannot resolve.

This does not mean the model is without value. It means the \emph{claim} that pretraining is
responsible for the performance is not supported by this experiment. Distinguishing the two
requires either much more statistical power, or datasets on which transfer is larger. Neither
is available to us here, and we suspect neither is available in the current fNIRS literature.

The finding also explains why the paper's own reported gains might be small and heterogeneous:
the model authors report cross-subject improvements of roughly $+5.0$, $+4.3$ and $+1.2$ pp on
different datasets, none individually significant. Our independent estimate on a clean dataset
is in the same range, which is at least consistent.

\subsection{Architecture, not pretraining}
The significant contrast in our results is the transformer representation against hand-crafted
features, and it survives with a random encoder. This is the random-features phenomenon
\citep{rahimi2007random} reappearing in a new modality: a sufficiently expressive fixed
representation can be competitive with a learned one when the downstream task is small. For
practice, this suggests that a well-designed encoder architecture may deliver most of the value
that foundation-model framing promises, at a fraction of the data-curation cost---and that the
field's claims should be framed accordingly.

\subsection{Systemic signal and the cost of cleaning}
Our short-separation ablation is the result we least expected. Physiological regression is
recommended, is well motivated, and helps elsewhere \citep{fischer2026density}. Here it reduced
accuracy, and---more tellingly---it removed precisely the part of the signal on which the
embedding arms were outperforming classical features. The most plausible account is that the
auditory conditions differ in global physiology (arousal, attention, global haemodynamics) in
addition to local cortical responses, and short channels capture that global component.

This has two implications. First, studies that use short-separation regression should report
decoding performance \emph{with and without} it; a large drop is informative about where the
decodable variance lives. Second, and more broadly, it reinforces the case for interpreting
fNIRS decoding results alongside explicit tests of non-neural confounds.

\subsection{Integration must be learned, not concatenated}
Our stacking result is a negative one, and we report it because the temptation to combine
representations is strong. At matched dimensionality, concatenation degraded performance and was
statistically indistinguishable from the weaker component alone. The redundancy between
hand-crafted features and transformer embeddings is high; adding a representation is not the same
as adding information. A combined model would need a learned routing mechanism, and it should be
evaluated against the same random-encoder control.

\subsection{Implications for evaluation practice}
We suggest three concrete practices for fNIRS foundation-model papers: (i) a random-encoder
baseline with multiple seeds; (ii) explicit contamination accounting---which evaluation datasets
overlap the pretraining corpus; (iii) reporting of preprocessing sensitivity, in particular
short-separation regression, because it can change both absolute accuracy and the ranking of
methods. Each is inexpensive and each materially changes what a reported number means.
